\documentclass{article}
\usepackage{amsmath,amssymb,amsthm,algorithm,algorithmic}
\usepackage{hyperref}
\newtheorem{theorem}{Theorem}
\newtheorem{lemma}{Lemma}
\newtheorem{assumption}{Assumption}
\newcommand{\prox}{\operatorname{prox}}
\begin{document}

\section*{Proximal Gradient Method (PGM)}

\section*{Mathematical Formulation}
We consider the composite optimization problem  

\[
\min_{x\in\mathbb{R}^{d}} \; \Phi(x)\;:=\; f(x)+g(x),
\tag{1}
\]

where  

\begin{itemize}
  \item $f:\mathbb{R}^{d}\!\to\!\mathbb{R}$ is convex and differentiable with $L$‑Lipschitz continuous gradient,
  \[
      \|\nabla f(u)-\nabla f(v)\|\le L\|u-v\|\quad\forall u,v\in\mathbb{R}^{d},
  \tag{2}
  \]
  \item $g:\mathbb{R}^{d}\!\to\!\mathbb{R}\cup\{+\infty\}$ is proper, closed, convex (possibly non‑smooth).
\end{itemize}

For a stepsize $\eta>0$, the proximal operator of $g$ is  

\[
\prox_{\eta g}(z)\;:=\;\arg\min_{y\in\mathbb{R}^{d}}
\Big\{ g(y)+\frac{1}{2\eta}\|y-z\|^{2}\Big\}.
\tag{3}
\]

The **Proximal Gradient update** is  

\[
x_{k+1}\;=\;\prox_{\eta g}\!\bigl(x_{k}-\eta \nabla f(x_{k})\bigr),
\qquad k=0,1,2,\dots .
\tag{4}
\]

\section*{Pseudocode}
\begin{algorithm}[H]
\caption{Proximal Gradient Method}
\begin{algorithmic}[1]
\STATE \textbf{Input:} stepsize $\eta\in(0,1/L]$, initial point $x_{0}\in\mathbb{R}^{d}$, max iter $K$.
\FOR{$k=0$ \textbf{to} $K-1$}
    \STATE $u_{k}\gets x_{k}-\eta\,\nabla f(x_{k})$ \hfill \# gradient step
    \STATE $x_{k+1}\gets \prox_{\eta g}(u_{k})$ \hfill \# proximal step
\ENDFOR
\STATE \textbf{Output:} $x_{K}$
\end{algorithmic}
\end{algorithm}

\section*{Dry Run Example}
Take $f(w)=(w-3)^{2}$, $g\equiv0$ (so $\prox_{\eta g}= \text{id}$), stepsize $\eta=0.1$, start $w_{0}=0$.

\[
\nabla f(w)=2(w-3).
\]

\begin{center}
\begin{tabular}{c|c|c|c}
Iteration $k$ & $w_{k}$ & $\nabla f(w_{k})$ & $w_{k+1}=w_{k}-\eta\nabla f(w_{k})$ \\ \hline
0 & $0$ & $-6$ & $0-0.1(-6)=0.6$ \\
1 & $0.6$ & $-4.8$ & $0.6-0.1(-4.8)=1.08$ \\
2 & $1.08$ & $-3.84$ & $1.08-0.1(-3.84)=1.464$ \\
\end{tabular}
\end{center}

Objective values: $\Phi(w_{0})=9$, $\Phi(w_{1})=(0.6-3)^{2}=5.76$, $\Phi(w_{2})=(1.08-3)^{2}=3.6864$, $\Phi(w_{3})=(1.464-3)^{2}=2.365$.
Thus the method decreases the objective each iteration.

\newpage

\section*{Assumptions}
\begin{assumption}[Convexity and Smoothness of $f$]
$f$ is convex and differentiable with $L$‑Lipschitz gradient as in (2). Moreover $f$ is bounded below.
\end{assumption}
\begin{assumption}[Convexity of $g$]
$g$ is proper, closed, convex; its proximal operator $\prox_{\eta g}$ is well defined for all $\eta>0$.
\end{assumption}
\begin{assumption}[Stepsize]
The constant stepsize satisfies $0<\eta\le \frac{1}{L}$.
\end{assumption}

\section*{Main Convergence Theorem}
\begin{theorem}[Sublinear convergence of PGM]\label{thm:main}
Let $\{x_{k}\}$ be generated by (4) under Assumptions 1–3 and denote $x^{\star}\in\arg\min\Phi$. Then for all $K\ge1$  

\[
\Phi(x_{K})-\Phi(x^{\star})\;\le\;\frac{\|x_{0}-x^{\star}\|^{2}}{2\eta K}.
\tag{5}
\]

Consequently $\Phi(x_{k})\to\Phi(x^{\star})$ and the rate is $O(1/k)$.
\end{theorem}

\section*{Theoretical Properties}
\begin{itemize}
\item \textbf{Stability:} The iteration is a non‑expansive mapping because the proximal operator is firmly non‑expansive and the gradient step is $L$‑Lipschitz; with $\eta\le1/L$ the composition is contractive in the squared norm sense.
\item \textbf{Hyperparameter sensitivity:} The stepsize bound $\eta\le1/L$ is tight; larger $\eta$ may violate (5) and cause divergence.
\item \textbf{Comparison:} When $g\equiv0$, (4) reduces to gradient descent with the classic $O(1/k)$ rate. When $f\equiv0$, (4) becomes the proximal point algorithm, which also enjoys $O(1/k)$ in function values.
\end{itemize}

\section*{Discussion}
The proof relies solely on convexity, Lipschitz continuity of $\nabla f$, and the non‑expansiveness (firm non‑expansiveness) of $\prox_{\eta g}$:
\[
\|\prox_{\eta g}(u)-\prox_{\eta g}(v)\|^{2}
\le \langle u-v,\;\prox_{\eta g}(u)-\prox_{\eta g}(v)\rangle .
\tag{6}
\]
No stochasticity or variance terms appear; thus the deterministic bound (5) holds for any initialization.

\newpage

\section*{Preliminaries (Definitions and Lemmas)}
\begin{lemma}[Descent Lemma for $L$‑smooth $f$]\label{lem:descent}
For any $u,v\in\mathbb{R}^{d}$,
\[
f(v)\le f(u)+\langle\nabla f(u),v-u\rangle+\frac{L}{2}\|v-u\|^{2}.
\tag{7}
\]
\end{lemma}
\begin{proof}
Standard consequence of $L$‑Lipschitz gradient; see e.g. \cite{nesterov2004intro}.
\end{proof}

\begin{lemma}[Firm non‑expansiveness of prox]\label{lem:prox}
For any $u,v\in\mathbb{R}^{d}$ and any $\eta>0$,
\[
\|\prox_{\eta g}(u)-\prox_{\eta g}(v)\|^{2}
\le \langle u-v,\;\prox_{\eta g}(u)-\prox_{\eta g}(v)\rangle .
\tag{8}
\]
\end{lemma}
\begin{proof}
By definition of proximal operator as the minimizer of a strongly convex function; see \cite{parikh2014proximal}.
\end{proof}

\section*{Main Proof (Step‑by‑Step Derivation)}
Let $x^{\star}\in\arg\min\Phi$. Define the auxiliary point  

\[
y_{k}=x_{k}-\eta\nabla f(x_{k}) .
\tag{9}
\]

Then the update (4) reads $x_{k+1}= \prox_{\eta g}(y_{k})$.

\bigskip

\noindent\textbf{[Step 1]} (Apply Lemma~\ref{lem:prox} with $u=y_{k}$, $v=x^{\star}$):
\[
\|x_{k+1}-x^{\star}\|^{2}
\le \langle y_{k}-x^{\star},\,x_{k+1}-x^{\star}\rangle .
\tag{10}
\]

\noindent\textbf{[Step 2]} (Insert $y_{k}=x_{k}-\eta\nabla f(x_{k})$):
\[
\|x_{k+1}-x^{\star}\|^{2}
\le \langle x_{k}-x^{\star},\,x_{k+1}-x^{\star}\rangle
-\eta\langle\nabla f(x_{k}),\,x_{k+1}-x^{\star}\rangle .
\tag{11}
\]

\noindent\textbf{[Step 3]} (Expand the first inner product using the identity 
$\langle a,b\rangle=\frac12(\|a\|^{2}+\|b\|^{2}-\|a-b\|^{2})$):
\[
\langle x_{k}-x^{\star},\,x_{k+1}-x^{\star}\rangle
=\frac12\bigl(\|x_{k}-x^{\star}\|^{2}+\|x_{k+1}-x^{\star}\|^{2}
-\|x_{k}-x_{k+1}\|^{2}\bigr).
\tag{12}
\]

\noindent\textbf{[Step 4]} (Plug (12) into (11) and cancel $\|x_{k+1}-x^{\star}\|^{2}$ on both sides):
\[
\frac12\|x_{k+1}-x^{\star}\|^{2}
\le \frac12\|x_{k}-x^{\star}\|^{2}
-\frac12\|x_{k}-x_{k+1}\|^{2}
-\eta\langle\nabla f(x_{k}),\,x_{k+1}-x^{\star}\rangle .
\tag{13}
\]

\noindent\textbf{[Step 5]} (Rearrange (13)):
\[
\|x_{k}-x^{\star}\|^{2}
-\|x_{k+1}-x^{\star}\|^{2}
\ge \|x_{k}-x_{k+1}\|^{2}
+2\eta\langle\nabla f(x_{k}),\,x_{k+1}-x^{\star}\rangle .
\tag{14}
\]

\noindent\textbf{[Step 6]} (Use convexity of $g$: by definition of proximal operator,
\[
x_{k+1}= \prox_{\eta g}(y_{k})
\;\Longleftrightarrow\;
\frac{1}{\eta}(y_{k}-x_{k+1})\in\partial g(x_{k+1}).
\tag{15}
\]
Hence for any $z$,
\[
g(z)\ge g(x_{k+1})+\Big\langle\frac{1}{\eta}(y_{k}-x_{k+1}),\,z-x_{k+1}\Big\rangle .
\tag{16}
\]
Take $z=x^{\star}$ and substitute $y_{k}=x_{k}-\eta\nabla f(x_{k})$:
\[
g(x^{\star})\ge g(x_{k+1})+\Big\langle \frac{x_{k}-\eta\nabla f(x_{k})-x_{k+1}}{\eta},\,x^{\star}-x_{k+1}\Big\rangle .
\tag{17}
\]
\]

\noindent\textbf{[Step 7]} (Rewrite (17) to isolate the gradient term):
\[
\langle\nabla f(x_{k}),\,x_{k+1}-x^{\star}\rangle
\ge \frac{1}{\eta}\langle x_{k}-x_{k+1},\,x_{k+1}-x^{\star}\rangle
+g(x_{k+1})-g(x^{\star}).
\tag{18}
\]

\noindent\textbf{[Step 8]} (Plug (18) into inequality (14)):
\[
\|x_{k}-x^{\star}\|^{2}
-\|x_{k+1}-x^{\star}\|^{2}
\ge \|x_{k}-x_{k+1}\|^{2}
+2\bigl\langle x_{k}-x_{k+1},\,x_{k+1}-x^{\star}\bigr\rangle
+2\eta\bigl(g(x_{k+1})-g(x^{\star})\bigr).
\tag{19}
\]

\noindent\textbf{[Step 9]} (Observe that
$\|x_{k}-x^{\star}\|^{2}
-\|x_{k+1}-x^{\star}\|^{2}
= \|x_{k}-x_{k+1}\|^{2}
+2\langle x_{k+1}-x^{\star},\,x_{k}-x_{k+1}\rangle$,
so the first two terms on the right‑hand side of (19) cancel exactly.)
Thus we obtain the clean inequality  

\[
\|x_{k}-x^{\star}\|^{2}
-\|x_{k+1}-x^{\star}\|^{2}
\ge 2\eta\bigl(g(x_{k+1})-g(x^{\star})\bigr).
\tag{20}
\]

\noindent\textbf{[Step 10]} (Apply Lemma~\ref{lem:descent} with $u=x_{k}$, $v=x_{k+1}$):
\[
f(x_{k+1})\le f(x_{k})+\langle\nabla f(x_{k}),x_{k+1}-x_{k}\rangle
+\frac{L}{2}\|x_{k+1}-x_{k}\|^{2}.
\tag{21}
\]

\noindent\textbf{[Step 11]} (From optimality condition (15) we have 
$\frac{1}{\eta}(x_{k}-\eta\nabla f(x_{k})-x_{k+1})\in\partial g(x_{k+1})$, therefore
\[
\langle\nabla f(x_{k}),x_{k+1}-x_{k}\rangle
=-\frac{1}{\eta}\|x_{k}-x_{k+1}\|^{2}
-\langle s_{k+1},x_{k+1}-x_{k}\rangle,
\tag{22}
\]
where $s_{k+1}\in\partial g(x_{k+1})$.
Since $g$ is convex,
$\langle s_{k+1},x_{k+1}-x_{k}\rangle\ge g(x_{k+1})-g(x_{k})$.
Hence
\[
\langle\nabla f(x_{k}),x_{k+1}-x_{k}\rangle
\le -\frac{1}{\eta}\|x_{k}-x_{k+1}\|^{2}
-(g(x_{k+1})-g(x_{k})).
\tag{23}
\]

\noindent\textbf{[Step 12]} (Insert (23) into (21)):
\[
f(x_{k+1})\le f(x_{k})
-\frac{1}{\eta}\|x_{k}-x_{k+1}\|^{2}
-(g(x_{k+1})-g(x_{k}))
+\frac{L}{2}\|x_{k+1}-x_{k}\|^{2}.
\tag{24}
\]

\noindent\textbf{[Step 13]} (Collect terms and use $\eta\le 1/L$ so that
$\frac{L}{2}-\frac{1}{\eta}\le -\frac{1}{2\eta}$):
\[
f(x_{k+1})+g(x_{k+1})
\le f(x_{k})+g(x_{k})
-\frac{1}{2\eta}\|x_{k}-x_{k+1}\|^{2}.
\tag{25}
\]

\noindent\textbf{[Step 14]} (Denote $\Phi(x)=f(x)+g(x)$. Inequality (25) yields a monotone decrease:
\[
\Phi(x_{k+1})\le\Phi(x_{k})-\frac{1}{2\eta}\|x_{k}-x_{k+1}\|^{2}.
\tag{26}
\]

\noindent\textbf{[Step 15]} (Combine (20) and (26). From (20):
\[
\|x_{k}-x^{\star}\|^{2}
-\|x_{k+1}-x^{\star}\|^{2}
\ge 2\eta\bigl(g(x_{k+1})-g(x^{\star})\bigr).
\tag{27}
\]
Add $2\eta\bigl(f(x_{k+1})-f(x^{\star})\bigr)$ to both sides and use (26):
\[
\|x_{k}-x^{\star}\|^{2}
-\|x_{k+1}-x^{\star}\|^{2}
\ge 2\eta\bigl(\Phi(x_{k+1})-\Phi(x^{\star})\bigr)
-\|x_{k}-x_{k+1}\|^{2}.
\tag{28}
\]

\noindent\textbf{[Step 16]} (Replace $-\|x_{k}-x_{k+1}\|^{2}$ using (26):
\[
-\|x_{k}-x_{k+1}\|^{2}
\ge -2\eta\bigl(\Phi(x_{k})-\Phi(x_{k+1})\bigr).
\tag{29}
\]

\noindent\textbf{[Step 17]} (Insert (29) into (28)):
\[
\|x_{k}-x^{\star}\|^{2}
-\|x_{k+1}-x^{\star}\|^{2}
\ge 2\eta\bigl(\Phi(x_{k+1})-\Phi(x^{\star})\bigr)
-2\eta\bigl(\Phi(x_{k})-\Phi(x_{k+1})\bigr).
\tag{30}
\]

\noindent\textbf{[Step 18]} (Simplify (30)):
\[
\|x_{k}-x^{\star}\|^{2}
-\|x_{k+1}-x^{\star}\|^{2}
\ge 2\eta\bigl(\Phi(x_{k})-\Phi(x^{\star})\bigr)
-4\eta\bigl(\Phi(x_{k})-\Phi(x_{k+1})\bigr).
\tag{31}
\]

\noindent\textbf{[Step 19]} (Discard the non‑negative term $4\eta(\Phi(x_{k})-\Phi(x_{k+1}))$ to obtain a cleaner bound)
\[
\|x_{k}-x^{\star}\|^{2}
-\|x_{k+1}-x^{\star}\|^{2}
\ge 2\eta\bigl(\Phi(x_{k})-\Phi(x^{\star})\bigr).
\tag{32}
\]

\noindent\textbf{[Step 20]} (Rearrange (32):
\[
\Phi(x_{k})-\Phi(x^{\star})
\le \frac{1}{2\eta}\bigl(\|x_{k}-x^{\star}\|^{2}
-\|x_{k+1}-x^{\star}\|^{2}\bigr).
\tag{33}
\]

\noindent\textbf{[Step 21]} (Sum (33) for $k=0,\dots,K-1$; telescoping occurs):
\[
\sum_{k=0}^{K-1}\bigl(\Phi(x_{k})-\Phi(x^{\star})\bigr)
\le \frac{1}{2\eta}\bigl(\|x_{0}-x^{\star}\|^{2}
-\|x_{K}-x^{\star}\|^{2}\bigr)
\le \frac{\|x_{0}-x^{\star}\|^{2}}{2\eta}.
\tag{34}
\]

\noindent\textbf{[Step 22]} (Since $\Phi(x_{k})$ is non‑increasing, the smallest term among the first $K$ iterates is $\Phi(x_{K})$. Hence
\[
K\bigl(\Phi(x_{K})-\Phi(x^{\star})\bigr)
\le \sum_{k=0}^{K-1}\bigl(\Phi(x_{k})-\Phi(x^{\star})\bigr).
\tag{35}
\]

\noindent\textbf{[Step 23]} (Combine (34) and (35)):
\[
\Phi(x_{K})-\Phi(x^{\star})
\le \frac{\|x_{0}-x^{\star}\|^{2}}{2\eta K}.
\tag{36}
\]
This is exactly the claim (5) of Theorem~\ref{thm:main}. $\square$

\section*{Rate Derivation (With Constants)}
From (36) we can read the explicit constant:  

\[
C:=\frac{\|x_{0}-x^{\star}\|^{2}}{2\eta},
\qquad\text{so that}\qquad
\Phi(x_{K})-\Phi(x^{\star})\le \frac{C}{K}.
\]

If the stepsize is chosen at the maximal admissible value $\eta=1/L$, then  

\[
C=\frac{L}{2}\,\|x_{0}-x^{\star}\|^{2},
\quad\text{and}\quad
\Phi(x_{K})-\Phi(x^{\star})\le \frac{L\|x_{0}-x^{\star}\|^{2}}{2K}.
\]

\section*{Special Cases}
\begin{itemize}
\item \textbf{Pure gradient descent} ($g\equiv0$): $\prox_{\eta g}=\mathrm{id}$, the proof reduces to the classical $O(1/k)$ bound for gradient descent on convex $L$‑smooth functions.
\item \textbf{Proximal point algorithm} ($f\equiv0$): the update becomes $x_{k+1}= \prox_{\eta g}(x_{k})$. The same algebra yields $\Phi(x_{K})- \Phi(x^{\star})\le \frac{\|x_{0}-x^{\star}\|^{2}}{2\eta K}$, i.e. $O(1/k)$ for the pure proximal method.
\item \textbf{Elastic net regularization} $g(x)=\lambda_{1}\|x\|_{1}+\frac{\lambda_{2}}{2}\|x\|^{2}$: the prox operator has a closed‑form soft‑thresholding plus shrinkage; the same $O(1/k)$ rate holds.
\end{itemize}

\begin{thebibliography}{9}
\bibitem{nesterov2004intro}
Y.~Nesterov,
\emph{Introductory Lectures on Convex Optimization: A Basic Course},
Kluwer Academic Publishers, 2004.

\bibitem{parikh2014proximal}
N.~Parikh and S.~Boyd,
\emph{Proximal Algorithms},
Foundations and Trends in Optimization, 1(3):123–231, 2014.
\end{thebibliography}

\end{document}